\chapter{Analiz Modeli}

\section{Genel Bakış}
Analiz modeli, sistemin iş mantığını ve kullanıcı etkileşimlerini anlamak için oluşturulmuştur. Bu model, sistem gereksinimlerini görselleştirmek ve sistemin davranışını tanımlamak için UML diyagramları kullanmaktadır.

\section{Use Case Diagram}
Use Case diyagramı, sistemin aktörleri (kullanıcılar) ve sistemle olan etkileşimlerini (use case'ler) gösterir.

\begin{figure}[H]
    \centering
    \includegraphics[width=\textwidth]{use-case.png}
    \caption{Use Case Diagram}
    \label{fig:use-case}
\end{figure}

\textbf{Aktörler:}
\begin{itemize}
    \item Authenticated User (Kimlik Doğrulanmış Kullanıcı)
\end{itemize}

\textbf{Ana Use Case'ler:}
\begin{enumerate}
    \item Kayıt Ol
    \item Giriş Yap
    \item Profil Yönetimi
    \item Birebir Sohbet Oluştur
    \item Grup Sohbeti Oluştur
    \item Sohbet Listesi Görüntüle
    \item Mesaj Gönder
    \item Mesaj Geçmişi Görüntüle
    \item Mesaj Sil
    \item Dosya Yükle
    \item Dosya İndir
    \item Bildirim Al
\end{enumerate}

\section{Domain Model (Etki Alanı Modeli)}
Domain model, sistemin temel iş varlıklarını (entities) ve aralarındaki ilişkileri gösterir.

\begin{figure}[H]
    \centering
    \includegraphics[width=\textwidth]{domain-model.png}
    \caption{Domain Model Diagram}
    \label{fig:domain-model}
\end{figure}

\textbf{Ana Varlıklar:}
\begin{itemize}
    \item \textbf{User (Kullanıcı):} id, email, username, passwordHash, firstName, lastName, profilePicture, statusMessage, createdAt, updatedAt
    \item \textbf{Chat (Sohbet):} id, type (DIRECT, GROUP), name, createdBy, createdAt, updatedAt
    \item \textbf{ChatParticipant (Sohbet Katılımcısı):} id, chatId, userId, role (MEMBER, ADMIN), joinedAt
    \item \textbf{Message (Mesaj):} id, chatId, senderId, content, messageType (TEXT, FILE, IMAGE, VIDEO), fileUrl, status (SENT, DELIVERED, READ, DELETED), createdAt, updatedAt
    \item \textbf{File (Dosya):} id, originalName, storedName, mimeType, size, url, uploadedBy, createdAt
    \item \textbf{Notification (Bildirim):} id, userId, type (MESSAGE, FILE, CHAT_INVITE), title, body, data, read, createdAt
\end{itemize}

\textbf{Varlık İlişkileri:}
\begin{itemize}
    \item User 1..* ChatParticipant
    \item Chat 1..* ChatParticipant
    \item Chat 1..* Message
    \item User 1..* Message
    \item User 1..* File
    \item User 1..* Notification
\end{itemize}

\section{Activity Diagram}
Activity diagram, sistem içindeki iş akışlarını gösterir.

\begin{figure}[H]
    \centering
    \includegraphics[width=0.8\textwidth]{activity-diagram.png}
    \caption{Mesaj Gönderme Activity Diagram}
    \label{fig:activity-message}
\end{figure}

\begin{figure}[H]
    \centering
    \includegraphics[width=0.8\textwidth]{activity-diagram-file-upload.png}
    \caption{Dosya Yükleme Activity Diagram}
    \label{fig:activity-file}
\end{figure}

Aşağıdaki ana akışlar modellenmiştir:
\begin{enumerate}
    \item \textbf{Mesaj Gönderme Akışı:} Mesaj gönderme sürecini gösterir.
    \item \textbf{Dosya Paylaşımı Akışı:} Dosya yükleme ve paylaşma sürecini gösterir.
\end{enumerate}

Her activity diagram, sürecin adımlarını, karar noktalarını ve paralel aktiviteleri gösterir.
